\section{Related Work}
\label{sec:related}

\paragraph{Training-Free Visual Anomaly Detection.}
A growing body of work tackles anomaly detection without task-specific training by leveraging powerful pretrained representations.
PatchCore~\citep{roth2022patchcore} introduced a coreset-based memory bank of patch-level features, enabling efficient nearest-neighbor anomaly scoring that remains competitive on industrial benchmarks.
WinCLIP~\citep{jeong2023winclip} extended this to the zero-shot setting using CLIP-based window-level similarity scores.
Self-supervised vision transformers---particularly DINOv2~\citep{oquab2024dinov2}---have provided rich patch features that further improve nearest-neighbor methods, and FoundAD~\citep{li2026foundad} showed that feature adaptation strategies can close the gap between training-free and fine-tuned approaches.
Despite strong performance on MVTec and VisA, these methods have been evaluated almost exclusively on industrial surfaces; their generalization across medical, retail, road, and logical anomaly domains remains unexplored.
AnomalyClaw incorporates patch-level experts as one component of a broader expert pool, rather than relying on a single detector as the sole decision-maker.

\paragraph{VLM-Based Anomaly Detection.}
Vision-language models offer semantic reasoning that purely visual detectors lack.
AnomalyGPT~\citep{gu2024anomalygpt} fine-tunes a VLM to produce anomaly judgments and localization maps.
Anomaly-OneVision~\citep{cao2025anomalyonevision} distills anomaly knowledge into a unified VLM handling multiple industrial categories.
The MMAD benchmark~\citep{jiang2024mmad} formalized VLM-based anomaly detection as visual question answering, revealing that even state-of-the-art VLMs struggle with fine-grained defect recognition.
AD-Copilot~\citep{zhang2026adcopilot} proposed an interactive assistant paradigm.
A common limitation is the absence of quantitative grounding and self-verification: VLMs make single-pass judgments with no mechanism to catch their own errors.
AnomalyClaw addresses both gaps through adversarial debate grounded by structured expert evidence.

\paragraph{Agent-Based Anomaly Detection.}
Tool-augmented agents have been applied to compensate for VLM perceptual weaknesses.
AgentIAD~\citep{zhang2025agentiad} equips a VLM with external tools for feature extraction and comparison, training the agent via supervised fine-tuning.
EAGLE~\citep{chen2026eagle} injects PatchCore-derived visual prompts (heatmaps, bounding boxes) directly into the VLM's visual input.
AutoIAD~\citep{wei2025autoiad} coordinates multiple specialized agents to automate the detection pipeline.
These systems demonstrate the value of combining VLMs with specialized detectors but use \emph{fixed} integration strategies: the order and type of tool invocation is predetermined regardless of query difficulty.
AnomalyClaw differs by making expert selection and reasoning depth \emph{adaptive}: the autonomous controller decides which experts to invoke and how many debate rounds to run based on the evolving state of the debate.
\Cref{tab:related_comparison} summarizes these differences.

\begin{table}[t]
\centering
\caption{%
    \textbf{Comparison with VLM-based methods for anomaly detection.}
    AnomalyClaw is the only system that (i)~uses adversarial debate for self-verification, (ii)~supports an extensible expert pool with autonomous invocation, (iii)~requires no training, and (iv)~is evaluated across 10 diverse domains.
}
\label{tab:related_comparison}
\small
\setlength{\tabcolsep}{2.8pt}
\begin{tabular}{@{}lcccccc@{}}
\toprule
Method & Expert & Channel & Self-Verify & Training & Adaptive & Cross-Domain \\
\midrule
AgentIAD~\citeyear{zhang2025agentiad}  & PatchCore+CLIP & Tool calls & No & SFT & No & No \\
EAGLE~\citeyear{chen2026eagle}         & PatchCore & Visual overlay & No & None & No & Partial \\
AutoIAD~\citeyear{wei2025autoiad}      & Multi-agent & Pipeline & No & None & No & No \\
DMAD~\citeyear{li2025dmad}             & None & Debate & Partial & None & No & No \\
\textbf{AnomalyClaw}  & \textbf{Extensible pool} & \textbf{Text context} & \textbf{Yes} & \textbf{None} & \textbf{Yes} & \textbf{Yes (10)} \\
\bottomrule
\end{tabular}
\end{table}

\paragraph{Multi-Agent Debate and Reasoning.}
Multi-agent debate---where multiple LLM instances critique and refine each other's answers---has shown promise for improving factual accuracy.
DMAD~\citep{li2025dmad} applied multi-agent deliberation to anomaly detection through iterative discussion among VLM agents.
M3MAD-Bench~\citep{liang2026m3mad} studied multi-agent debate for multimodal tasks, identifying ``collective delusion'' as a key failure mode where agents reinforce rather than correct errors.
VipAct~\citep{li2026vipact} proposed a vision-perception-action framework with role-based collaboration.
These systems use \emph{symmetric} debate where both agents play identical roles.
AnomalyClaw uses \emph{asymmetric} debate with structurally different roles (Proposer vs.\ Advocate), which avoids collective delusion by design: the Advocate's sole objective is to challenge anomaly claims, preventing the echo-chamber effect observed in symmetric debate~\citep{liang2026m3mad}.
Furthermore, AnomalyClaw grounds the debate with quantitative expert evidence rather than relying solely on VLM perception.

\paragraph{Cross-Domain Anomaly Detection Benchmarks.}
MMAD~\citep{jiang2024mmad} covers industrial defect detection across MVTec-AD and VisA but does not extend beyond manufactured objects.
ADNet~\citep{li2025adnet} and M3-AD~\citep{chen2025m3ad} broaden coverage to additional industrial categories.
BMAD~\citep{bao2024bmad} addresses medical domains but does not consider non-medical settings.
Our evaluation framework curates tasks from 10 domains under one unified protocol, enabling a broader study of how VLM-based detectors generalize across fundamentally different notions of visual anomaly.
